\documentclass[11pt]{article}

\usepackage[margin=1in]{geometry}
\usepackage{graphicx}
\usepackage{caption}
\usepackage{float}
\usepackage{amsmath}
\usepackage[colorlinks=true, linkcolor=blue, urlcolor=blue]{hyperref}

\captionsetup{font=small, labelfont=bf, justification=raggedright,
              singlelinecheck=false}


\graphicspath{{plots/}{./}}


\newcommand{\figab}[1]{%
  \IfFileExists{plots/#1_paper.pdf}%
    {\includegraphics[width=\textwidth]{#1_paper.pdf}}%
    {\IfFileExists{#1_paper.pdf}%
      {\includegraphics[width=\textwidth]{#1_paper.pdf}}%
      {\includegraphics[width=\textwidth]{#1.pdf}}}}

\begin{document}

\begin{center}
{\Large\bfseries Plots and Explanations}\\[6pt]
{\normalsize August 21, 2026}
\end{center}

\bigskip

\noindent
This note presents the descriptive results behind our research question: when
tensions between the United States and China rise, do US institutional
investors reallocate away from firms whose supply chains run through China?
Figures~\ref{fig:roll} through~\ref{fig:vw} track US holdings of European
firms sorted by the firms' own exposure to China. Figures~\ref{fig:cagg}
through~\ref{fig:cew} repeat the sort at the country level, using the three
exposure measures defined below. Figures~\ref{fig:b2} and~\ref{fig:b3} turn
the question around and ask whether firms held by US investors cut their China
ties pre-emptively. An important distinction:
a gap that opens between exposed and unexposed firms can reflect investors
trading, or it can reflect the set of firms classified as exposed changing
over time. The figures are paired so that the two explanations can be told
apart.

Here are the findings. Figure~\ref{fig:roll} appears to show US investors
moving heavily into China-linked European firms over two decades, and out of
firms without those ties. That appearance is an artefact of how the data were
collected: the set of firms counted as China-linked grows roughly sixfold as
the relationship records fill in. Compare each firm only with itself and the
two groups end up in essentially the same place (Figure~\ref{fig:fixed}). At
the country level there is no single answer: whether high-exposure countries
gained or lost US portfolio share depends on which of the three exposure
measures does the ranking, and two of them point in opposite directions
(Section 3). Firms held by US investors, meanwhile, did not cut their China
ties as tensions rose. Their supply-chain links to China grew about twice as
fast as those of firms with no US owners (Section 4).

\section*{1. Data and measurement}

\textbf{Holdings.} Institutional equity holdings come from FactSet Ownership.
The sample is European firms held by US-domiciled institutions, quarterly. The
holdings data begin in 1999 Quarter 1; the supply-chain relationship data that classify
firms begin later, so each figure is plotted over the full span its own
construction supports, and the
spans therefore differ across figures. For every fund we keep its last
complete report inside each calendar quarter and stamp it to quarter end. Dollar values are deflated by the
US CPI and expressed in 2020 dollars. When we report a group of firms' share
of the US portfolio, the denominator is the value of all equity positions
held by the same US institutions worldwide, not only their European
positions, so the share reads as the weight of that group in the investors'
overall book.

\textbf{Supply-chain exposure.} Firm-to-firm relationships come from FactSet
Revere, which records customer and supplier links between companies. We
define a firm's exposure to China as
\begin{equation*}
s_{i,t} \;=\;
\frac{\text{active links of firm } i \text{ to Chinese customers or suppliers in quarter } t}
     {\text{active links of firm } i \text{ to customers or suppliers anywhere in quarter } t},
\end{equation*}
counting only relationships recorded as active. A firm covered by the
relationship data but with no active link has $s_{i,t}=0$; a firm outside the
relationship data has no exposure value and is left unclassified rather than
treated as a zero.

Here, we report two kinds of exposure.  On the extensive margin, a firm either has a Chinese counterparty or it does
not. So we can measure whether a firm is entangled with China at
all.  On the intensive margin, a firm that already has Chinese
counterparties may have more or fewer of them, so we can see the variation.  We therefore report both rather than choosing one:
Figures~\ref{fig:roll} and~\ref{fig:fixed} sort firms on the extensive margin,
Figures~\ref{fig:ew} and~\ref{fig:vw} sort them on the intensive margin among
firms that have any Chinese counterparty, and all three country-level measures
below are continuous in $s_{i,t}$ and so carry both.

\textbf{Country-level exposure.} We aggregate firm exposure to the country
level in three ways, which differ in whose links count and how firms are
weighted. The \emph{aggregate} measure divides the total number of China
links of all firms in a country by the total number of their links to any
counterparty, so firms with many relationships dominate. The
\emph{value-weighted} measure averages the firm-level shares $s_{i,t}$ using
market capitalization as weights. It is computable only where market value is
observed, and that coverage is neither flat nor monotone: about 71 percent of
firm-quarters in 2003, a peak 90 percent in 2007, then a decline to 43
percent by 2023. The measure tilts toward larger listed firms. The
\emph{equal-weighted} measure averages $s_{i,t}$ with equal weights, so small firms count as much as large ones. The three measures answer slightly different questions, and we show all three because, as
Section 3 documents, the answer to ``did US investors leave high-exposure
countries'' turns out to depend on which one is used.

\textbf{Tension.} The background series in every figure is the USA$|$China
direction of a bilateral geopolitical-risk index built from newspaper text, in
which the United States is the acting party and China the counterparty\footnote{The quarterly value is the level of
the index in the closing month of each quarter. It is from \url{https://www.matteoiacoviello.com/ai_gpr.html}}. Each figure carries the index
twice: once as a top panel with its own axis, and once as a grey silhouette
behind the data panels. The index is low through 2017, apart from an
isolated spike at the end of 2016 (8.2, against a 1999--2017 median of 0.9). It
jumps to a new level with the 2018--2019 tariff rounds. It eases after that but
stays well above its pre-2018 average through 2020--2021, falls back in 2022, and spikes again in 2023.

\textbf{The reference date.} Two constructions in this note are
anchored at 2017 Quarter 4, and the choice of that quarter is deliberate. The first US tariffs of the trade war arrive during 2018. The quarter ending 31 December 2017 is therefore the last clean quarter before the trade war.

\textbf{The bottom panel of every figure.} Every figure carries a bottom panel
reporting how many units lie behind each line, and the full history is
plotted rather than a window we picked. Two artefacts show up there. First,
the relationship data arrive in batches. In the rolling firm figures, the
count of firms with no China link climbs from about 430 in 2003 to roughly
6,600 by 2012, in discrete loads: jumps of 53 percent in 2007Q3, 96 percent in
2008Q4, 48 percent in 2010Q4 and 53 percent in 2011Q1. Only after 2011Q1 does
the count grow smoothly. Second, many US funds filed holdings semiannually in
the early sample. Through 2004Q2, about 1,720 funds report in the first and
third quarters against 4,620 in the second and fourth, with identical
positions per filing (285 against 290), so the difference is who reported, not
what was reported. Fund counts converge in 2004Q3, but odd-quarter filings
stay thinner for several years (201 positions per fund against 280 in 2004,
310 against 324 in 2006) and only match by 2008--2009. Quarter-to-quarter
readings before 2009 are therefore uninformative, even where the annual
pattern is sound. Neither artefact is a reason to hide the early data. Both
are reasons to read each line against the counts below it.

\textbf{Timing.} Whenever firms or countries are sorted into groups, the sort
uses information dated $t-1$ and outcomes are measured at $t$, so group
membership is always predetermined relative to the holdings it is paired
with. Figure~\ref{fig:fixed} applies the same rule at both ends of each
quarter: a firm contributes to a quarter's growth only if it sits in the same
group at the start and the end of that quarter.

\section*{2. US holdings sorted by firm-level exposure}

Figure~\ref{fig:roll} splits European firms each quarter into those with no
active China link and those with at least one, and follows US institutional
holdings of the two groups, from 2003Q3, when firms can first be classified, to
2023Q4. Read naively, the figure tells a dramatic story.
Holdings of firms with China ties climb from under \$20 billion at the start to
\$1.4 trillion in 2023, an index of 145 against a 2017 base of 100. Holdings of
firms without ties move the other way, from a peak of \$1.01 trillion at the end
of 2013 down to \$654 billion, an index of 73. The portfolio shares tell the same
story. Measured against the US institutions' worldwide equity book, the tied group
rises from under 1 to 8 percent while the untied group falls from a peak near 10
percent to under 4 percent. The two lines cross around 2015. Taken at face value this would say US investors tilted
\emph{toward} China-linked firms across two decades of rising tensions.

The two groups, however, are not the same firms over time. In 2012 the
China-linked group contains 145 firms; by 2023 it contains 858. The newcomers
are largely not firms that began trading with China in the meantime. They are
firms whose existing relationships entered the data late: Revere's coverage
grows in the discrete loads documented above, and when a firm's links are
finally recorded, the firm jumps from the no-ties line to the China-ties line.
No money is traded when it jumps. Both lines move because the rosters change.

Figure~\ref{fig:fixed} takes the jumps out. Each quarter, growth in each group
is computed only over firms that are present in that quarter and the one
before, and that sit in the same group at both dates. The quarterly growth
rates are then chained into an index. A firm that enters the data, leaves it,
or switches groups contributes nothing in the quarter it moves, so roster
changes cannot move these lines by construction, and no reference date is
involved in who gets compared; the 2017Q4 = 100 base only fixes where the two
lines cross 100. The rotation does not survive. The chained index ends 2023 at
105 for firms without China ties and 96 for firms with them, on matched
samples of 8,445 and 769 firms. Of the 72-point gap in Figure~\ref{fig:roll},
at most ten points can be about trading. The rest is the rosters, and we
attach no interpretation to the ten points.

Figures~\ref{fig:roll} and~\ref{fig:fixed} sort on the extensive margin alone,
so by construction they are silent on whether investors respond to the depth of
a firm's China exposure rather than its existence.
Figures~\ref{fig:ew} and~\ref{fig:vw} supply that second margin. They drop the
unexposed firms as a sorting dimension. Instead they split the firms that do have
Chinese counterparties into quartiles of $s_{i,t-1}$, with cutpoints computed
equal-weighted in Figure~\ref{fig:ew} and value-weighted in Figure~\ref{fig:vw},
mirroring the two firm-level aggregation schemes of Section 1. The unexposed group
is carried along as a reference line only. If US investors were shedding
the most exposed names, the top-quartile line should sag after 2018. It does the
opposite. Indexed to 2017, top-quartile holdings reach 252 in the equal-weighted
sort and 188 in the value-weighted sort. The bases are small: those groups end the
sample at \$122 billion and \$396 billion, 0.7 and 2.3 percent of the worldwide
book.

The two sorts disagree about whether the response is monotone in exposure, which
is why the middle quartiles are reported numerically in the captions even though
the figures draw only the extremes.
Under the value-weighted cutpoints the end-of-sample index rises with exposure
throughout: 102, 143, 164, 188 from the bottom quartile to the top. Under the
equal-weighted cutpoints it does not: 52, 220, 282, 252, so the third quartile
outgrows the fourth. We would not read a dose response into either pattern. What
both sorts do agree on is the absence of the predicted one, since no quartile
retreats after 2018. These are also rolling sorts, so reclassification
contributes to all four lines.

\begin{figure}[p]
  \centering
  \figab{fig_A_firm_zero_vs_positive}
  \caption{\textbf{US holdings of European firms with and without China
    ties.} This figure plots US institutional holdings of European firms,
    split each quarter into firms with no active Chinese customer or supplier
    relationship at the end of the previous quarter and firms with at least
    one. The top panel shows the bilateral US--China tension index. The grey
    silhouette in the lower panels repeats it. The lower panels give, for each
    group, the real value of US institutional holdings (billions of 2020 dollars,
    log scale), the same series indexed to 100 at end-2017, and the group's share
    of all equity held by the same US institutions worldwide. End-2017 is the last
    quarter before the 2018 tariff rounds. Firms not covered by the relationship
    data in the prior quarter are unclassified, so the groups need not sum to the
    total.
    Two cautions. Membership is reassigned every quarter, and the count of
    China-linked firms grows from 145 in 2012 to 858 in 2023 as the relationship
    data deepen, so the widening gap partly reflects composition;
    Figure~\ref{fig:fixed} removes it. The split is also on the
    extensive margin alone. A firm that halves its Chinese suppliers, or replaces
    one with a supplier elsewhere, stays in the same group here.
    Figures~\ref{fig:ew} and~\ref{fig:vw} sort on the depth of exposure instead.}
  \label{fig:roll}
\end{figure}

\begin{figure}[p]
  \centering
  \figab{fig_A_firm_zero_vs_positive_chained}
  \caption{\textbf{The same comparison with the roster changes removed.} This
    figure re-draws the contrast of Figure~\ref{fig:roll} as a chain-linked
    matched-sample index. For each quarter, growth in each group is computed
    only over firms present in both that quarter and the previous one and
    classified in the same group at both dates; the quarterly growth rates are
    then chained, and both lines are drawn at 100 in 2017Q4. That base is
    presentation, not construction: it fixes where the lines cross 100, not
    which firms are compared. A firm whose links enter the data, or that
    switches groups, contributes nothing in the quarter it moves, so the
    roster growth visible in Figure~\ref{fig:roll} cannot move these lines.
    The bottom panel counts the matched firms behind each quarter's growth,
    8,445 and 769 in 2023Q4. The 72-point rotation of Figure~\ref{fig:roll}
    does not survive: the two groups end within ten points of each other, 105
    against 96, and we attach no interpretation to the residual. The zigzag
    before 2009 is the odd-quarter reporting gap described in Section 1, and
    it thins the matched counts in those quarters too.}
  \label{fig:fixed}
\end{figure}

\begin{figure}[p]
  \centering
  \figab{fig_A_firm_quartile_ew}
  \caption{\textbf{US holdings by quartile of firm exposure, equal-weighted
    sort.} This figure sorts European firms with a positive China link share into
    quartiles of that share at the end of the previous quarter. Quartile cutpoints
    give each firm equal weight. Plotted are US institutional holdings of
    the lowest and the highest quartile, with the no-ties group as a reference.
    Panels are as in Figure~\ref{fig:roll}:
    tension index, real holdings (log scale), holdings indexed to 100 at
    end-2017, and share of the US institutions' worldwide equity book. The
    no-ties group is drawn separately rather than folded into the bottom
    quartile because it is not the low tail of the same distribution. A bottom-quartile firm has a
    China share that is measured and merely small. The no-ties line pools two
    different things: firms with a positive supply-chain denominator and no Chinese
    counterparty (4,276 firms in 2023Q4), and firms with no supply-chain links at
    all (4,241, of which 2,926 have no recorded relationship of any type). The second kind is the larger part of that line
    for most of the sample: pooled across quarters, 65 percent of no-ties
    firm-quarters have no supply-chain link to measure. The two middle quartiles are
    computed but not drawn, because four sequential blues on a log axis are not
    separable by eye. Their end-of-sample index values are 220 for the second
    quartile and 282 for the third, against 52 for the first and 252 for the
    fourth, so the response is not monotone in exposure under these cutpoints:
    the third quartile ends above the fourth.}
  \label{fig:ew}
\end{figure}

\begin{figure}[p]
  \centering
  \figab{fig_A_firm_quartile_mcapw}
  \caption{\textbf{US holdings by quartile of firm exposure, value-weighted
    sort.} As Figure~\ref{fig:ew}, but quartile cutpoints weight firms by
    market capitalization, and the whole sort, including the no-ties
    reference group, is restricted to firms with observed market value at the
    end of the previous quarter so that all lines describe one universe. This
    restriction matters for interpretation, and coverage is neither flat nor
    monotone: it runs near 70 percent of firm-quarters in 2003, rises to about
    90 percent by 2007, then falls to 44 percent by the end of the sample. The
    figure therefore conditions on coverage and tilts toward larger listed
    firms. Under these cutpoints the end-of-sample index rises with exposure
    throughout, 102, 143, 164 and 188 from the first quartile to the fourth,
    which the equal-weighted sort of Figure~\ref{fig:ew} does not.}
  \label{fig:vw}
\end{figure}

\section*{3. US holdings sorted by country-level exposure}

Figures~\ref{fig:cagg}, \ref{fig:cvw}, and~\ref{fig:cew} move the sort to the
country level. That is the unit of the companion country-by-quarter panel of
28 countries; the main triple difference is estimated one level down, on firm
by holder-group by quarter cells. Each quarter, countries are
ranked on their exposure to China at the end of the previous quarter, using
in turn the aggregate, value-weighted, and equal-weighted measures of
Section 1, and split at the cross-country 25th and 75th percentiles into a
bottom quartile, middle half, and top quartile. A country enters the ranking
only in quarters where at least twenty of its firms can be classified.

Two mechanical warnings apply to all three figures. First, the bottom-quartile
group is economically tiny, about \$9 billion and five basis points of the
worldwide book at the end of the sample, so its line carries little
information. Second, group membership turns over at a low and
roughly constant rate throughout: under the aggregate measure an
average of 2.5 rankable countries change groups per quarter through 2016 and
2.3 per quarter afterward, and the other two rankings are just as flat. What
makes the early bottom-quartile line unreliable is not the volume of turnover
but its leverage. That group holds only six or seven countries, so one large
country crossing the 25th-percentile cut moves the line by hundreds of index
points: when Spain displaces Greece in 2015Q2 the group's index goes from 119
to 687 and back to 103 the next quarter. The early-sample swings are that, and
should not be read as trading. The aggregate measure has a related fragility
we have not yet fixed: its denominator is a link count with no floor, so a
country with hundreds of firms but only a handful of recorded links can post
an extreme exposure value. Poland enters the top quartile in four consecutive
quarters of 2010 and 2011 on two recorded links. How to screen such cells is
an open item, and the ranking below applies no link floor.

The substantive result is that the three measures disagree, and the
disagreement is itself the finding. Indexed to 100 at end-2017, holdings in
top-quartile countries stand in 2023 at 141 under the aggregate ranking, 99
under the value-weighted ranking, and 84 under the equal-weighted ranking.
Under the aggregate measure, which is driven by the link counts of the
largest, most connected firms, high-exposure countries also gain portfolio
share (from about 4 percent in 2012 to 7.2 percent in 2023). Under the
equal-weighted measure, which lets many small firms set a country's rank,
high-exposure countries lose ground relative to the middle. A conclusion
about whether US investors left high-exposure \emph{countries} therefore
depends on how exposure is aggregated, which is why the regressions that
follow this note use all three measures side by side rather than committing
to one, and why no single descriptive figure should be read as the answer.

\begin{figure}[p]
  \centering
  \figab{fig_A_country_m1_quartiles}
  \caption{\textbf{US holdings by country exposure quartile, aggregate
    measure.} This figure ranks countries each quarter on the aggregate
    exposure measure at the end of the previous quarter, that is, total China
    links of the country's firms over their total links to any counterparty. It
    then splits them at the cross-country 25th and 75th percentiles. Panels show
    the tension index, then for each group real US institutional holdings on a log
    scale, holdings indexed to 100 at end-2017, and the group's share of the US
    institutions' worldwide equity book. The log scale is necessary: the two
    extreme groups differ by two orders of magnitude in size. Countries with fewer than twenty
    classifiable firms in a quarter are excluded that quarter. Group
    membership turns over at a similar rate throughout, two to three countries
    per quarter, but the bottom quartile holds only six or seven countries, so
    the large early swings of its line reflect single countries crossing a
    percentile boundary, not trading.}
  \label{fig:cagg}
\end{figure}

\begin{figure}[p]
  \centering
  \figab{fig_A_country_m2_quartiles}
  \caption{\textbf{US holdings by country exposure quartile, value-weighted
    measure.} As Figure~\ref{fig:cagg}, but countries are ranked on the
    value-weighted measure: the average of firm-level China link shares
    weighted by market capitalization. Because the measure requires market
    value, it is computed on the covered subset of firms only, which shrinks
    over time and skews toward larger listed companies; the ranking should be
    read jointly with the aggregate and equal-weighted versions rather than
    on its own.}
  \label{fig:cvw}
\end{figure}

\begin{figure}[p]
  \centering
  \figab{fig_A_country_m3_quartiles}
  \caption{\textbf{US holdings by country exposure quartile, equal-weighted
    measure.} As Figure~\ref{fig:cagg}, but countries are ranked on the
    equal-weighted measure: the simple average of firm-level China link
    shares, one firm one vote. Firms with supply-chain links but none to China
    enter this average as zeros, and they are not invisible to the aggregate
    measure of Figure~\ref{fig:cagg} either: they supply about half its
    denominator, 49 percent at the end of 2023 and 52 percent pooled over the
    sample, so they pull that measure down rather than passing through it. The
    firms invisible to the aggregate measure are those covered by the
    relationship data with no active supply-chain link at all, which enter
    neither its numerator nor its denominator; they are 44 percent of
    classified firms at the end of 2023 and 61 percent pooled. The divergence
    between the two rankings reflects that asymmetry together with how
    concentrated China links are in a handful of highly connected firms.}
  \label{fig:cew}
\end{figure}

\section*{4. Do firms with US investors cut their China links?}

Government pressure could also work through firms rather than investors. A
European firm that values its US shareholder base might cut its China ties
pre-emptively, so as not to lose it. Figures~\ref{fig:b2} and~\ref{fig:b3}
test for this by sorting firms on their US institutional ownership at the end
of the previous quarter and following the growth of their China links. To
keep the comparison honest, two rules apply. Growth is computed on a cohort fixed
at $t-1$ and required to have measured link counts a year earlier and in the
current quarter, so expanding data coverage cannot manufacture growth. Links are
summed within each group before growth is taken, which keeps firms with zero links
in the denominator.

The hypothesis fails in a clean and instructive way. Firms with US investors add China links
\emph{faster} than firms without them through most of the tension era. That holds
in 18 of the 24 quarters from 2018 on, in every quarter of 2023, with an average
year-over-year gap of about ten percentage points. The six exceptions are scattered
and mostly small. The largest is a 5.7-point reversal in 2019Q4. At the end of the sample, year-over-year link
growth is 24
percent for US-held firms against 13 percent for firms with no US
institutional ownership, and in the three-way split of Figure~\ref{fig:b3}
the gradient is monotone, 24.5 percent for firms above the median US
holding, 21.8 percent below the median, 12.7 percent with none. Intensity
tells the same story: US-held firms above the median end the sample with 2.6
China links per firm, against 0.23 below the median and 0.09 for firms with
no US owners. One panel needs a warning. The total-links index shows the no-US group growing
fastest from 2017 on. That group averages fewer than one link per ten firms, so a
tiny base drives the line, not de-linking by US-held firms.
The most US-held firms show no de-linking at any point after 2017: their
year-over-year link growth never turns negative and their total China links
rise from 1,245 to 3,411. The one line that falls is the below-median US-held
group, whose links peak at 371 in 2020Q4 and are down to 257 by 2023Q1 before
recovering to 302, still 29 percent above their 2017Q4 level. That bucket is
re-sorted at each quarter's median, so firms crossing the median drain links
out of it; the two US-held arms taken together never decline.

\begin{figure}[p]
  \centering
  \figab{fig_B_link_growth_us2}
  \caption{\textbf{Growth of China links, firms with versus without US
    investors.} This figure sorts European firms on US institutional
    ownership at the end of the previous quarter into firms with none and
    firms with any, and plots the growth of their supply-chain links to
    China. Five panels: the tension index; year-over-year growth in each group's
    total China links; total links indexed to 100 at end-2017; China links per
    firm; and the number of cohort firms in each group, on a log scale. Growth is
    computed on a cohort fixed at the sort date and required to have measured links
    a year earlier and currently, so expanding data coverage cannot generate
    growth. Links are summed within the group before growth is taken, which keeps
    firms with zero links in the denominator. The ownership grid is zero-filled, so every firm has a row in
    every quarter and ``no US investors'' is a coded zero rather than a blank.
    Coverage of the underlying filings is a separate matter: through 2008Q2 the
    odd-quarter reporting gap described in Section~1 inflates the no-US bucket
    in the June and December quarters, which sort on a March or September
    holding, and 17 to 30 percent of that bucket held US money in both adjacent
    quarters. From 2008Q3 the same check flags 0.1 percent, so across the
    tension era ``no US investors'' is a measured zero and not missing
    coverage.}
  \label{fig:b2}
\end{figure}

\begin{figure}[p]
  \centering
  \figab{fig_B_link_growth_us3}
  \caption{\textbf{Growth of China links by intensity of US ownership.} As
    Figure~\ref{fig:b2}, but firms with US investors are further split at
    each quarter's median US dollar holding into a below-median and an
    above-median group. The split uses the dollar value of the position
    rather than the ownership share. Market value, which the share requires, is
    unobserved for a growing part of this figure's cohort as a whole, covering
    89 percent of firm-quarters in 2007 but only 43 percent in 2023, and its
    availability is size-correlated. That constraint does not bind on the firms
    actually being split: among firm-quarters with a positive US position,
    market value at the end of the previous quarter is observed for 98 percent,
    so a share-based sort would have been feasible. What the dollar sort costs
    is that it loads heavily on firm size, with a within-quarter rank
    correlation near 0.9 between position value and market capitalization; a
    share-based split would move about 13 percent of firm-quarters across the
    median. We accept and note the trade-off. Link growth is fastest for the
    most US-held firms at the end of the sample and from mid-2020 on, though
    the below-median group grew faster through 2017 and 2018; nothing here
    looks like pre-emptive de-linking.}
  \label{fig:b3}
\end{figure}

\section*{5. Taking stock}

Three descriptive facts come out of the nine figures. First, the apparent
decade-long rotation of US portfolios toward China-linked European firms is
composition: compare each firm only with itself and the two groups end the
sample within a few percent of each other, a residual too small and too
sensitive to construction for us to interpret. Second, at the country level the
sign of the reallocation depends on how exposure is aggregated, so the
question needs a regression design with all three measures, not a chart with
one. Third, firms with US owners did not cut China ties as tensions rose;
they added them faster than firms without US owners. None of this is
evidence that US institutions ignored the tensions; it says that whatever
response exists is small relative to the trends in the data and will have to
be identified off sharper variation than raw levels, which is what the
panel design this note accompanies is built to do.

\end{document}
